\section{The original simultaneous loss estimate}
\label{sec:global-loss-v035}

\status{LFS1--LFS17 prove the literal
$\texttt{LossFS}(\mathcal F,0)$ from the same original
\texttt{ChosenTT}, \texttt{ChosenBadp}, \texttt{LinkG}, and
\texttt{ChosenWS}.  The proof includes both the prime-logarithmic
\texttt{IOsum} bound and the universally read \texttt{ASK}-minus-
\texttt{Ga} bound.  It preserves the chosen integer functions and
their realizing member.  Existence of the compatible family and its
local gain estimates remains an upstream input.}

Throughout the final application, $\mathcal F:\texttt{FamP}(\epsilon)$
is the original supplied family and $0<\epsilon<1/52$.  Write
$D=\mathcal F.D(x,i)$ and $L=\mathcal F.L(x,i)$ only after fixing the
common threshold.  To keep the source's two small parameters distinct,
write
\[
 z=x^\epsilon,\quad u=x^{\epsilon/3},\quad
 b=e^z,\quad a=xb,\quad X=e^u=X',\quad
 b'=xb/X,\quad w''=b/(nX),\quad X''=xnX.                 \tag{LFS.1}
\]
The relative top-window width is denoted $h$; the exceptional-share
cutoff is $\theta=x^{-\epsilon}$.  Thus $h$ is the source field
\texttt{Par.th}, whereas $\theta$ is \texttt{Par.theta}.
All finite sums below are over the displayed sets; repeated indices
never denote independently chosen arithmetic data.

% BEGIN LFS1
\paragraph{LFS1: elementary prime-logarithmic estimates with their constants.}
\label{claim:LFS1}
For every real $Y\ge1$ and every finite set of primes $S\subset[1,Y]$,
\[
 \sum_{p\in S}\log p\le (\log4)Y,\qquad
 \sum_{p\in S}\frac{\log p}{p}\le\log Y+\log4.          \tag{LFS.2}
\]
Here is a proof of the exact constant needed later.  Let
$P(N)=\prod_{p\le N}p$.  Strong induction gives $P(N)\le4^N$.
The cases $N\le2$ are direct.  If $N=2r\ge4$, $N$ is composite and
$P(N)=P(N-1)$.  If $N=2r+1\ge3$, the product of the primes in
$(r+1,2r+1]$ divides $\binom{2r+1}{r}$: each occurs once in the
numerator factorial and not in either denominator factorial.
The two equal central binomial coefficients are among terms summing
to $2^{2r+1}$, so $\binom{2r+1}{r}\le4^r$.  Hence
$P(N)\le P(r+1)4^r\le4^{2r+1}$.  Taking logarithms and then
$N=\lfloor Y\rfloor$ proves the first estimate.

For the second, each $p^{\lfloor N/p\rfloor}$ divides $N!$, and
these factors for different primes are coprime.  Consequently
\[
 \sum_{p\le N}\lfloor N/p\rfloor\log p\le\log(N!)
       \le N\log N.
\]
Multiply the desired prime sum by $N$, use
$N/p\le\lfloor N/p\rfloor+1$, and apply the first estimate.
Division by $N>0$ proves (LFS.2) at integers; restriction to
$\lfloor Y\rfloor$ gives the real version.

We will also use, for an integer $p_0\ge3$ and any finite
$S\subset\{p_0,p_0+1,\ldots\}$,
\[
 \sum_{n\in S}\frac{\log n}{n^2}
 \le \frac{2(\log(p_0-1)+1)}{p_0-1}.                    \tag{LFS.3}
\]
Indeed $f(t)=\log t/t^2$ is decreasing for $t\ge2$, because
$f'(t)=(1-2\log t)/t^3\le0$; the inequality
$\log2\ge1/2$ follows by integrating $1/t\ge1/2$ on $[1,2]$.
Thus the full tail is at most
$\int_{p_0-1}^\infty f(t)\,dt=(\log(p_0-1)+1)/(p_0-1)$,
which is stronger than (LFS.3).  These arguments establish all three
estimates without a prime number theorem or a distribution theorem
for residue classes.
% END LFS1

% BEGIN LFS2
\paragraph{LFS2: the dyadic small-order prime bound.}
\label{claim:LFS2}
Let $e\ge1$, $J\ge1$, and $Y\ge1$.  For an odd prime $p$, let
$\bar p$ be its class in
$(\mathbb Z/2^{e+2}\mathbb Z)^\times/\{\pm1\}$.
Then every finite set $S$ of such primes with $p\le Y$ and
$\operatorname{ord}(\bar p)\le J$ satisfies
\[
 \sum_{p\in S}\frac{\log p}{p}
 \le\frac{12J\log^2(2^eY)}{2^e}.                       \tag{LFS.4}
\]
We give the residue calculation, rather than treating the order
condition as a prime-density assertion.  If $v$ is odd, one of
$v-1,v+1$ has 2-adic valuation exactly one.  Factoring
$v^2-1=(v-1)(v+1)$ therefore gives, for $r\ge1$,
\[
 2^{r+1}\mid v^2-1
 \quad\Longleftrightarrow\quad
 2^r\mid v-1\ \hbox{or}\ 2^r\mid v+1.
\]
Also $v^2+1\equiv2\pmod4$.  Induction on $j$, using these two
facts at each squaring, gives for $0\le j\le e$
\[
 v^{2^j}\equiv\pm1\pmod{2^{e+2}}
 \quad\Longleftrightarrow\quad
 v\equiv\pm1\pmod{2^{e+2-j}}.                          \tag{LFS.5}
\]
In particular the quotient group has exponent dividing $2^e$.
An element of order at most $J$ thus has order $2^r$, where
$r\le j_0:=\min(e,\lfloor\log_2J\rfloor)$.  Equation (LFS.5)
puts $p$ in one of the two classes $\pm1$ modulo
$M=2^{e+2-j_0}$.  We have $M\ge4$ and $2^e\le JM$.

Every integer $p\ge3$ in these two classes is $qM+1$ or $qM-1$
with $q\ge1$.  If $p\le Y$, then $q\le\lfloor Y\rfloor$ and
$p\ge qM/2$.  Extending the nonnegative sum along the two
progressions and using
$\sum_{q=1}^Nq^{-1}\le1+\log N$ proves
\[
 \sum_{p\in S}\frac{\log p}{p}
 \le\frac{4\log Y(1+\log Y)}M.
\]
Put $\lambda=\log(2^eY)$.  Since $e\log2\ge1/2$,
$\log Y(1+\log Y)\le3\lambda^2$.
The inequality $M^{-1}\le J/2^e$ now proves (LFS.4).
Thus the source's universal order constant may be taken to be
$C_{\rm ord}=12$.
% END LFS2

% BEGIN LFS3
\paragraph{LFS3: first-kind primes, including the narrow power bands.}
\label{claim:LFS3}
Let $X\ge2$, $0\le\rho\le1/2$, and $p_0\ge8$.
If a finite set of integers lies in
$[P(1-\rho),P(1+\rho)]\cap[p_0,\infty)$, with $1\le P\le X$,
then its cardinality is at most $2\rho P+1$.  Each reciprocal is
at most $\min(2/P,1/p_0)$ and each logarithm at most $\log(2P)$.
It follows that
\[
 \sum \frac{\log p}{p}
 \le (4\rho+1/p_0)\log(2X).                            \tag{LFS.6}
\]
For $p\ge p_0$ and $i\ge1$, the condition
$X(1-\rho)\le p^i\le X(1+\rho)$ implies
\[
 X^{1/i}(1-\rho)\le p\le X^{1/i}(1+\rho).
\]
Moreover $i\log8\le\log(2X)\le2\log X$, and
$\log8\ge2$.  One proof of the last inequality is
$\log2=2\int_0^{1/3}(1-t^2)^{-1}dt\ge2/3$.
Therefore $i\le\lfloor\log X\rfloor$.
Sum (LFS.6) over those exponents.  Overlaps can only increase this
upper bound.

Combining with LFS2 proves the following exact first-class bound.
If each prime of $S$ either has dyadic quotient order at most $J$
and lies below $Y$, or lies in one of the displayed bands, then
\[
 \sum_{p\in S}\frac{\log p}{p}
 \le\frac{12J\log^2(2^eY)}{2^e}
  +\lfloor\log X\rfloor(4\rho+1/p_0)\log(2X).          \tag{LFS.7}
\]
For the source parameters use
$J=\lfloor x^{\epsilon/4}\rfloor$,
$Y=\max(P.Y,X'')$, and $X=P.X'$.
An integer order strictly below $x^{\epsilon/4}$ is at most $J$.
Since $\log X'=x^{\epsilon/3}$, the number of bands is exactly
the source's $I_b=\lfloor x^{\epsilon/3}\rfloor$.
% END LFS3

% BEGIN LFS4
\paragraph{LFS4: the finite requested sums and the high-prime boundary.}
\label{claim:LFS4}
For a finite set $I\subset\mathbb N$ of exponents, write
\[
 A(I,A_0,p,w)=\sum_{r\in I}\max(A_0/p^r-w_r,0).
\]
If $p\ge2$, $0\le A_0\le3a/2$, all $r\in I$ are odd and
$w_r\ge0$, then the geometric series over positive odd exponents
gives
\[
 A(I,A_0,p,w)\le A_0\frac{p^{-1}}{1-p^{-2}}
 \le\frac{2a}{p},\qquad
 2\log p\,A(I,A_0,p,w)\le4a\frac{\log p}{p}.           \tag{LFS.8}
\]
For $d\ge4$, all positive exponents instead give
$A(I,A_0,d,w)\le A_0/(d-1)\le2a/d$ and hence the
discriminant contribution is at most $2a\log d/d$.

The sharper high-prime estimate uses the actual positive lower
weight $w''$, with $X''w''=a$, $X''\ge2$, and
$A_0\le a(1+h)$, $0\le h\le1$.
For $p>X''$, every term with exponent at least two is zero:
$p^2\ge2X''$ and $A_0/p^r\le w''$.  The exponent-one term
also vanishes for $p\ge X''(1+h)$; in the remaining interval it
is at most $w''h$.  There are at most $X''h+1$ integers there,
and $\log p\le\log(2X'')$.  Therefore
\[
 \sum_{p>X''}2\log p\,A(I_p,A_0,p,w_p)
 \le2\log(2X'')\bigl(ah^2+w''h\bigr).                 \tag{LFS.9}
\]
The endpoint allowance $+1$ is retained in the $w''h$ term.
All these estimates apply to finite subsets without any assumption
that every exponent or every prime is present.
% END LFS4

% BEGIN LFS5
\paragraph{LFS5: exact deficit splitting into four prime classes.}
\label{claim:LFS5}
Fix finite sets of lifts $J$, tops $T$, and primes $\mathcal P$;
put $m=|J|$ and $q=|T|$.  For each prime write $A_{j,t}(p)$ for
the finite requested sum in LFS4.  Let $g_j(p)$ be its retained
gain, $\tau(p)$ its correction, and $T_p\subset T$ its retained
tops.  The finite-place unit rule says that all places over $p$
have these same quantities and their normalized weights sum to
$\log p$, with zero contribution off $\mathcal P$.  Hence
\[
 \mathrm{ASK}-\mathrm{Ga}
 =\sum_{p\in\mathcal P}2\log p\left[
       \sum_{j,t}A_{j,t}(p)-\left(\sum_jg_j(p)-\tau(p)\right)
                                 \right]+\sum_{j,t}D_{j,t}.    \tag{LFS.10}
\]
This follows by interchanging finite sums; $D_{j,t}$ denotes the
discriminant term and the factor $2$ occurs on both sides.

At a served prime suppose the actual ledger inequality is
\[
 \sum_j\sum_{t\in T_p}(A_{j,t}(p)-\ell(p))-mf(p)
        \le\sum_jg_j(p)-\tau(p).
\]
Subtract and split $T=T_p\sqcup(T\setminus T_p)$.  The bracket
in (LFS.10) is at most
\[
 \sum_j\sum_{t\in T\setminus T_p}A_{j,t}(p)
                  +m|T_p|\ell(p)+mf(p).                       \tag{LFS.11}
\]
At an unserved prime the stipulated gain and correction are zero,
so the bracket is the entire requested sum.  Partition
$\mathcal P$ into served, first, second, and high primes, multiply
these inequalities by $2\log p\ge0$, and add them.  This proves
the source's deficit splitting with no positivity assumption on
the net gain and no change in the sets $T_p$.
% END LFS5

% BEGIN LFS6
\paragraph{LFS6: the four served-prime error bounds.}
\label{claim:LFS6}
Assume $C,f,\mathcal S,Q_1,Q_2,\Lambda\ge0$, $a\ge1$,
$k,K_0>0$, $p_0\ge3$, and $x,X,X''\ge1$.
The original envelopes are
\begin{align*}
 F(p)&=CfN^2\log_p a+Cf(\log_pa)^2,\\
 C_p&=C\mathcal S\log_p(2a),\\
 P_p&=\begin{cases}Q_1&p<p_0,\\Q_2&p\ge p_0,\end{cases}\\
 O_p&=C\log x\left(\frac{\Lambda a}{pk}
  +\begin{cases}ap^{-L_1(p)}&p<p_0,\\a/p^2&p\ge p_0.\end{cases}\right).
\end{align*}
Suppose served primes are at most $X''$, the $P_p,O_p$ terms
occur only for $p<X$, and $K_0\le p^{L_1(p)-1}$ for $p<p_0$.
Then their sums with the weight $2\log p$ satisfy
\begin{align*}
 \sum 2\log p\,F(p)
 &\le2CfX''\bigl(N^2\log a+2(\log a)^2\bigr),\\
 \sum 2\log p\,mqC_p
 &\le2C\mathcal S\log(2a)\,mqX'',\\
 \sum_{p<X}2\log p\,P_p
 &\le2\log4\,(p_0Q_1+XQ_2),\\
 \sum_{p<X}2\log p\,O_p
 &\le2C\log x\,a\left[
    (\Lambda/k+1/K_0)(\log X+\log4)
    +\frac{2(\log(p_0-1)+1)}{p_0-1}\right].             \tag{LFS.12}
\end{align*}
For the first line cancel $\log p$ in $\log_pa$, and use
$1/\log p\le2$ in the square term.  The number of served
primes is at most $\lfloor X''\rfloor\le X''$.
The second line is exact cancellation of $\log p$ followed by
the same cardinality bound.  For the third line, bound the low
prime sum by $\log4\,p_0$ and the whole sum below $X$ by
$\log4\,X$, using LFS1; counting the high contribution also
at low primes only increases this upper bound.
Finally $p^{-L_1(p)}\le1/(K_0p)$ below $p_0$.
Apply the reciprocal-prime bound in LFS1 to both terms with
$1/p$, and its square tail bound to the terms with $1/p^2$.
These steps prove every line of (LFS.12) with the stated constants.
% END LFS6

% BEGIN LFS7
\paragraph{LFS7: the complete eleven-term loss decomposition.}
\label{claim:LFS7}
Use the exact $\texttt{Par}$ quantities
\begin{equation}
\begin{aligned}
 f&=2n,\qquad\mathcal S=S_c+k+\eta,\qquad
 Q_1=Cr_Gb'\log x,\\
 Q_2&=C\left(b'/x+\eta+S_c+(w''+t)\log x/\delta_1\right)+b_o,
 \qquad\delta_1=x^{-\epsilon/8}.
\end{aligned}                                             \tag{LFS.13}
\end{equation}
Assume the eleven clauses of the finite ledger:
the unit rule in LFS5; $0\le a_t\le a(1+h)$;
odd prime exponents with weights at least $w''$;
positive discriminant exponents with nonnegative weights;
served primes at most $X''$; high primes greater than $X''$;
the served ledger in LFS5; zero unserved gains and corrections;
$f(p)\le F(p)$; $\ell(p)\le C_p+\mathbf1_{p<X}(P_p+O_p)$;
and $K_0\le p^{L_1(p)-1}$ below $p_0$.
Assume also the first-prime alternative of LFS3.

Take $h\le1/2$, so $a_t\le3a/2$.  LFS5, LFS4, and LFS6
give the following terms.  The first-class contribution is split
by LFS3 into its order and band parts.  At a served prime, the
envelope of $\ell(p)$ is nonnegative, so
$m|T_p|\ell(p)\le mq(C_p+\mathbf1_{p<X}(P_p+O_p))$ even
if $\ell(p)$ is negative.  Thus
\begin{equation}
\begin{aligned}
 \mathrm{ASK}-\mathrm{Ga}
 &\le E_F+E_C+E_{P,<}+E_{P,\ge}+E_O+E_{\rm out}\\
 &\qquad{}+E_{\rm ord}+E_{\rm band}+E_H+E_d+E_{\rm second}.
\end{aligned}                                               \tag{LFS.14}
\end{equation}
where the eleven terms, in order, are
\begin{align*}
 E_F&=m\,2CfX''(N^2\log a+2\log^2a),\\
 E_C&=2C\mathcal S\log(2a)mqX'',\\
 E_{P,<}&=mq\,2\log4\,p_0Q_1,\\
 E_{P,\ge}&=mq\,2\log4\,XQ_2,\\
 E_O&=mq\,2C\log x\,a
 \left[(\Lambda/k+1/K_0)(\log X+\log4)
       +\frac{2(\log(p_0-1)+1)}{p_0-1}\right],\\
 E_{\rm out}&=\sum_{p\ {\rm served}}4ma\frac{\log p}{p}|T\setminus T_p|,\\
 E_{\rm ord}&=4mqa\frac{C_{\rm ord}J_o
                       \log^2(2^{e_2}\max(Y,X''))}{2^{e_2}},\\
 E_{\rm band}&=4mqaI_b(4\rho+1/p_0)\log(2X),\\
 E_H&=mq\,2\log(2X'')(ah^2+w''h),\\
 E_d&=mq\,2a\log d/d,\\
 E_{\rm second}&=4mqa\sum_{p\ {\rm second}}\frac{\log p}{p}.
\end{align*}
Here $J_o=\lfloor x^{\epsilon/4}\rfloor$,
$I_b=\lfloor x^{\epsilon/3}\rfloor$, and $C_{\rm ord}=12$.
The terms correspond literally to
\texttt{tFloors}, \texttt{tCp}, \texttt{tPrLow},
\texttt{tPrHigh}, \texttt{tOr}, the omitted-top sum,
\texttt{tOrders}, \texttt{tBand}, \texttt{tHigh},
\texttt{tD}, and the second-class sum.  No first-class or
discriminant term has been silently absorbed into a different
arithmetic object.
% END LFS7

% BEGIN LFS8
\paragraph{LFS8: the same small-share inequality pays three errors.}
\label{claim:LFS8}
Let $\mathcal C$ be the candidate set and $\Sigma_p\ge0$.
Suppose the original chosen-weight inequality is
\[
 W:=\sum_{p\in\mathcal C}\frac{\log p}{p}
    \left(2\mathbf1_{\{\Sigma_p>\theta\}}+4\Sigma_p\right)
       \le x^{-9}.                                      \tag{LFS.15}
\]
Suppose every second-class prime belongs to $\mathcal C$ and has
$\Sigma_p>\theta$.  Then
\[
 E_{\rm second}\le2amqW\le2amq x^{-9}.                  \tag{LFS.16}
\]
Indeed the indicator part of $2amqW$ is exactly the required
$4amq\log p/p$, and its remaining summands are nonnegative.
If, for each candidate, the original error satisfies
$io(p)\le2\Sigma_p(2mqa/p)$, then the $4\Sigma_p$ part of
$amqW$ proves simultaneously
\[
 \sum_{p\in\mathcal C}\log p\,io(p)
       \le amqW\le amq x^{-9}.                          \tag{LFS.17}
\]

Let $\mathcal R$ be the original large-order residue-prime set.
Suppose an omitted served top is bad at that prime, and the prime
lies in $\mathcal R$.  Inclusion of these top sets and nonnegativity
of their weights show
\[
 E_{\rm out}\le4am\sum_{p\in\mathcal R}\frac{\log p}{p}
                         |\{t\in T:\mathrm{bad}(p,t)\}|.
\]
The original chosen-top inequality bounds this by
\[
                      E_{\rm out}\le4amq x^{-2+\epsilon}.
                                                               \tag{LFS.18}
\]
If a served prime is absent from $\mathcal R$, the omission
condition says $T\setminus T_p$ is empty, so this extension of
the sum is valid also at those primes.  Equations
(LFS.16)--(LFS.18) use precisely the same $\Sigma$, bad predicate,
and degree functions; a second independent choice is unnecessary.
% END LFS8

% BEGIN LFS9
\paragraph{LFS9: uniform absorption with the original leading constant and thirteen errors.}
\label{claim:LFS9}
Fix $0<\epsilon<1/4$, $e\ge0$, $K\ge1$ and
$C_{\rm ord}=12\le K$.  Assume the original
\texttt{WindowsD} inequalities:
\begin{gather*}
 x^2\le2^{e_2}\le n\le x^e,\qquad x^2\le k\le x^e,
 \qquad 0\le N,t,r_G\le x^e,\\
 x^4+1\le p_0\le x^e,\quad0\le\Lambda\le x^{2\epsilon},
 \quad0\le S_c,\eta\le e^{x^{\epsilon/4}},\\
 0\le\rho\le x^{-2},\quad0\le h\le x^{-1},\quad
 x^2\le K_0,d,\quad0\le C\le K,\quad
 1\le Y\le2X,\quad0\le b_o\le\sqrt b.                 \tag{LFS.19}
\end{gather*}
There is one threshold depending only on $\epsilon,e,K$ such that
every such parameter record and every finite ledger in LFS7--LFS8
satisfy, with $U=qmb$ and $\sigma=100Kx^{4\epsilon}/x$,
\[
 \mathrm{ASK}-\mathrm{Ga}
       \le(2\log4\,C+13\sigma)U,
 \qquad \sum_{p\in\mathcal C}\log p\,io(p)\le\sigma U,
 \qquad0\le\sigma\le1.                                \tag{LFS.20}
\]
Here $q\ge1$ and $m\ge0$.  We prove the uniform size step
explicitly; we do not infer it from the later, different
\texttt{Win\_at} predicate.

Put $E=x^e$, $S=e^{x^{\epsilon/4}}$, $A=\log x+z$, and
\[
 L_1=\log(2xEX),\qquad L_2=\log(2xE^2X).
\]
For sufficiently large $x$, substitution of (LFS.19) into the
eleven displayed terms gives the following upper coefficients
after factoring out $U$.  For $E_{P,\ge}$ the listed expression
is the coefficient after subtracting $2\log4\,C$.
\begin{center}\small
\begin{tabular}{p{0.20\textwidth}p{0.72\textwidth}}
Term & Upper coefficient\\\hline
$E_F$ & $4KxE^2(X/b)(E^2A+2A^2)$\\
$E_C$ & $2K(2S+E)(A+\log2)xE(X/b)$\\
$E_{P,<}$ & $2\log4\,KE^2x\log x/X$\\
High precision
 & $\begin{aligned}
 &2\log4\{K[2XS/b\\
 &\quad{}+(x^{-2}+XE/b)x^{\epsilon/8}\log x]+X/\sqrt b\}
 \end{aligned}$\\
$E_O$ & $\begin{aligned}
 &2Kx\log x[(z^2+1)x^{-2}(u+\log4)\\
 &\qquad{}+2(e\log x+1)x^{-4}]
 \end{aligned}$\\
$E_{\rm out}$ & $4x^{-1+\epsilon}$\\
$E_{\rm ord}$ & $4Kx^{-1+\epsilon/4}L_2^2$\\
$E_{\rm band}$ & $4xu(4x^{-2}+x^{-4})(u+\log2)$\\
$E_H$ & $2L_1(x^{-1}+x^{-3}/X)$\\
$E_d$ & $4\log x/x$\\
$E_{\rm second}$ & $2x^{-8}$\\
Prime-logarithmic $io$ sum & $x^{-8}$
\end{tabular}
\end{center}
For the floor term $1/q\le1$ was used.  The high-precision
term is worth checking separately: from $b'=xb/X$ and
$w''=b/(nX)$,
\[
 \frac{XQ_2}{b}
 =C\left[1+\frac{X(\eta+S_c)}b
       +\left(\frac1n+\frac{Xt}b\right)
                         x^{\epsilon/8}\log x\right]
        +\frac{Xb_o}b.                                    \tag{LFS.21}
\]
This is the unique leading term in the table.
For $E_O$ use $p_0-1\ge x^4$ and
$\log(p_0-1)\le e\log x$.  For the order term use
$2^{e_2}\max(Y,X'')\le2xE^2X$.  For the discriminant term,
$\log t/t$ is decreasing for $t\ge e^1$ and $d\ge x^2$,
so $2x\log d/d\le4\log x/x$; no upper bound on $d$ is
assumed.  These justify every row of the table.

Every coefficient in the table is $o(\sigma)$, uniformly in
all the parameter records satisfying (LFS.19).
For completeness the only non-exponential comparisons needed
are as follows.  Since $L_1,L_2=O_{e,\epsilon}(u)$, the ratios
to $\sigma$ of the order, band, and high rows are respectively
\[
 O(x^{-37\epsilon/12}),\qquad
 O(x^{-10\epsilon/3}),\qquad O(x^{-11\epsilon/3}).
\]
The principal part of the orbit row has ratio
$O(\log x\,x^{-5\epsilon/3})$; its other part is smaller.
The omitted-top row has ratio $O(x^{-3\epsilon})$.
The high-precision remainder without an exponential has ratio
$O(\log x\,x^{-1-31\epsilon/8})$.
The discriminant, second-class and $io$ rows have ratios
$O(\log x/x^{4\epsilon})$ and $O(x^{-7-4\epsilon})$.
All remaining ratios are polynomial-logarithmic factors times
one of $X/b$, $XS/b$, $X/\sqrt b$, or $1/X$, and tend to zero:
their logarithms tend to $-\infty$ because
$x^{\epsilon/3}+x^{\epsilon/4}=o(x^\epsilon)$ and
$\log x=o(x^r)$ for $r>0$.
These last limits follow, for example, by integrating
$1/t\le t^{r/2-1}$ and then dividing by $x^r$.

Choose a common threshold making each of the ten non-leading
loss terms at most $\sigma U$, the high-precision remainder at
most $3\sigma U$, and the $io$ sum at most $\sigma U$.
There are ten terms besides $E_{P,\ge}$ in (LFS.14), giving
exactly $10+3=13$.  Enlarge the threshold until
$100Kx^{4\epsilon-1}\le1$.  This proves (LFS.20), including
$m=0$ by multiplication of the coefficient inequalities rather
than division by $U$.

If desired the same threshold can be enlarged to satisfy all
three quantitative \texttt{LargeX} comparisons:
$100Kx^{4\epsilon}\le x$,
$100Kx^{4e+4}\le e^u$, and
$100Kx^{4e+4}e^{u+x^{\epsilon/4}}\le e^{z/2}$.
The proof uses a possibly larger cutoff than the source's
particular sufficient one, while retaining its exact
\texttt{joint\_core} constants and quantifier order.
% END LFS9

% BEGIN LFS10
\paragraph{LFS10: every original loss-window clause, uniformly on the family.}
\label{claim:LFS10}
Fix the constant $C$ and function $\delta(x)\to0$ supplied by
\texttt{LinkG}.  At each member put $b_o=b^{\delta(x)}$ and take
the literal parameter record \texttt{Loss.par}, with
$C_b=0$ and $C_{\rm ord}=12$.
The complete coordinate identification is
\[
\begin{array}{c|c@{\qquad}c|c}
 P.x&x&P.\epsilon&\epsilon\\
 P.e_2&D.R.e&P.n& D.R.n\\
 P.k&D.R.k&P.N& k+m(t+1)\\
 P.t&D.R.t&P.p_0&\lceil k^2x^{26}\rceil_+\\
 P.\Lambda&L.\mathrm{LaG}&P.r_G&r_G(\epsilon)\\
 P.S_c&L.\mathrm{Sc}=4kD.s.\mathrm{Del}&P.\eta&L.\mathrm{etaSz}\\
 P.\rho&G.\mathrm{rhoB}(D,\mathcal F.GK,L.sc)&P.h&1/(kx^{13})\\
 P.K_0&k^2x^{78}&P.d&D.host.d\\
 P.C&\max(C,0)&P.C_{\rm ord}&12\\
 P.Y&X'(1+\mathcal F.C_r\rho)&P.b_o&b^{\delta(x)}\\
 P.\beta&0&P.s_0&m|TT|x^{-12}.
\end{array}                                                   \tag{LFS.22}
\]
Choose
\[
                 e=\max(40,\mathcal F.C_t+3),\qquad
                 K=1+|C|+12.                                \tag{LFS.23}
\]
One threshold, independent of $i$, makes the entire original
\texttt{Windows}(P,e,K) true, and makes $k^2x^{26}\le X'$.
In particular its projection \texttt{WindowsD} is true.

Here are all the checks.  The original sizes in SZ1--SZ8 give,
for $x\ge3$,
\[
 n\le x^3,\quad m\le x,\quad x^4\le k\le x^5,
 \quad t\le x^{\mathcal F.C_t+\epsilon/100}
                   \le x^{e-2}.                              \tag{LFS.24}
\]
The source also has the sharper stored $t$-window; the weaker
bound displayed here suffices.  Its rank frame has
$2^{e_2}\le n\le o2^{e_2}$ and $o\le x^{\epsilon/100}$.
The lower $n$-window is at least $x^{2+\epsilon/100}$;
divide by this bound for $o$ to obtain $x^2\le2^{e_2}$.
Also $N=k+m(t+1)\le3x^{e-1}\le x^e$.
The ceiling in (LFS.22) satisfies
\[
 x^4+1\le p_0\le x^{36}+1\le x^{40}\le x^e.
\]
The lower inequality follows already from $k\ge x^4$ and
$x^{26}\ge2$; the upper one uses $k\le x^5$ and retains the
ceiling's $+1$.

The remaining scalar checks are
\begin{align*}
 0\le\Lambda&\le x^{\epsilon+\epsilon/100}\le x^{2\epsilon},\\
 S_c=4kD.s.\mathrm{Del}
 &\le4x^5e^{x^{\epsilon/4}/2}\le e^{x^{\epsilon/4}},
 &0<\eta&\le e^{x^{\epsilon/4}},\\
 0\le\rho&\le x^{-2},
 &0<h&=1/(kx^{13})\le x^{-1},\\
 x^2&\le k^2x^{78},
 &x^2&\le d.                                                \tag{LFS.25}
\end{align*}
For $S_c$ use \texttt{SizeDel}'s
$\log(D.s.\mathrm{Del})\le x^{\epsilon/4}/2$, and absorb
$4x^5$ into the other half exponential.  The $\eta$ inequality
is exactly its logarithmic size clause.  Nonnegativity of
$\rho$ follows by expanding its original expression into
nonnegative terms, using $\mathcal F.GK.\mathrm{Nonneg}$;
its upper bound is the literal \texttt{FamP.rho\_le} field.
For $d$ use \texttt{SizeQD}'s
$\log d\ge x^{19\mathcal F.ez/100}$.
Here $\mathcal F.ez=\epsilon/\mathcal F.J>0$ is fixed before
the threshold, so its exponential eventually exceeds $x^2$.
The integer $r_G(\epsilon)$ is fixed and is at most $x\le x^e$
eventually.

Since $\mathcal F.C_r\ge1$ is fixed, take $x\ge\mathcal F.C_r$.
Then $0\le\mathcal F.C_r\rho\le x^{-1}\le1$, proving
$1\le Y\le2X'$.  Since eventually $\delta(x)\le1/2$,
$0<b^{\delta(x)}\le e^{z/2}$; no sign assumption on
$\delta(x)$ is needed.  Clearly $0\le P.C\le K$,
$0\le P.C_{\rm ord}\le K$, and $0\le P.\beta=0\le Kx^{-2}$.
The top bound gives $|TT|\le2n\delta_1\le2x^3\le x^4$,
so $0\le s_0\le x^{-7}\le x^{-2}$.  These include the two
clauses discarded by \texttt{WindowsD}; that projection is
equivalently the full window for the record with
$\beta=s_0=0$.
Finally $k^2x^{26}\le x^{36}\le e^{x^{\epsilon/3}}$ eventually.
Every threshold used here depends only on the original fixed
family constants, $C$, and the eventual bound for $\delta$.
It precedes $x$ and every $i\in\mathcal F.I_x$.
% END LFS10

% BEGIN LFS11
\paragraph{LFS11: exact lift, top, gain, and prime-logarithmic dictionaries.}
\label{claim:LFS11}
Fix one original $D,L$ and $X_G:G.\mathrm{Inst}(D,L)$.
Let $J=\{1,\ldots,m\}$, so $|J|=m$, and let
$T=\mathrm{TTv}(D)=\{\operatorname{val}(t):t\in TT\}$.
The frame's integer-label map is injective and
$\mathrm{ofNat}(\operatorname{val}(t))=t$.  Thus
\[
 |T|=|TT|,\quad
 \sum_{r\in T}F(r)=\sum_{t\in TT}F(\operatorname{val}(t)).     \tag{LFS.26}
\]
The same bijection restricts to the bad subsets defined by
\[
 \mathrm{badAt}(p,r)=\neg G.\mathrm{Good}
 (X_G,p,X_G.bs(G.\mathrm{stepP}(p),\mathrm{ofNat}(r))).
\]
Their cardinalities are exactly \texttt{badTops}, not merely
bounded by those cardinalities.

Define the original finite prime universe
\[
 B=1+\max\bigl(\lfloor a+a/(kx^{13})\rfloor_+,
                                  \lfloor X''\rfloor_+\bigr),
 \qquad\mathcal P=\{p<B:p\text{ prime and inert in }K_d\}.    \tag{LFS.27}
\]
The candidate set is literally
$\mathcal C=\texttt{cands}(L,\mathcal F.GK)
 =\texttt{sAt}(L,\mathcal F.GK)$: each candidate has $p<X'$,
so the condition $p<\lceil X'\rceil_+$ adds nothing.
The source's rational-prime ledger \texttt{ldAt} has
$W=\mathcal P$, $\mathrm{below}=\mathrm{id}$ and weight
$\log p$.  Its unit rule is therefore a singleton-sum identity.
Its gain is
\[
 g_j(p)=\begin{cases}0&p\in L.\mathrm{drop},\\
 G.\mathrm{gainP}(X_G,\mathcal F.GK,a_L,b_L,p,j)&p\notin L.\mathrm{drop},
 \end{cases}
\]
and its correction is the original $G.\mathrm{tauP}$.

Let $\ell_w=\log N\mathfrak p_w/[E_0:\mathbb Q]$ for the
original coefficient field $E_0=D.\mathrm{what}_2.LL$, and
$e_w=v_w(p)$ for $w\mid p$.  The weight called
\texttt{Lnk.lE} is $\ell_we_w$.
The exact finite-place degree formula of LC3 gives
\[
                       \sum_{w\mid p}\ell_we_w=\log p.       \tag{LFS.28}
\]
The family support contains the full fiber over each $p\le X''$.
The original \texttt{ParametersFill} identities say, on that
support, that the gains and corrections are precisely the
rational-prime functions above.  At non-inert primes they
vanish by $L.\mathrm{not\_inert}$.  At $p>X''$, candidates
are absent, so $VV$ is empty, $\tau_P=io_P=0$, and the
explicit high-prime branch of $\mathrm{gainP}$ is zero.
At any further support places the summand is therefore zero.
Grouping finite sums by the full fibers in (LFS.28) proves
\[
 \boxed{\quad
 (\mathrm{ldAt}) .\mathrm{Ga}(J)=L.\mathrm{Ga},\qquad
 \sum_{p\in\mathcal C}\log p\,io_P(p)=\mathrm{IOsum}(L).
 \quad}                                                       \tag{LFS.29}
\]
The first identity uses
$\mathrm{sz}(w,y)=2\ell_we_wy$, hence retains the factor $2$.
The second uses the same $L.io(w)=io_P(p)$ from
\texttt{ParametersFill}; no factor $2$ is inserted into it.
Neither identity replaces the support by an infinite sum of
possibly nonzero terms.
% END LFS11

% BEGIN LFS12
\paragraph{LFS12: the actual four prime classes satisfy the order and omission alternatives.}
\label{claim:LFS12}
The source classifies a prime as served if $p\notin L.\mathrm{drop}$;
otherwise as high if $p>X''$; otherwise as second if it is a
candidate with $\Sigma_p>\theta$; otherwise as first.
We check these branches on the original universe (LFS.27).

First no type prime $T_u$ is inert.  The original norm witness
$\pi\in\mathcal O_E$ has $\pi\bar\pi=T_u$ and
$\pi\equiv1\pmod{M\mathcal O_E}$, where
$M=Q_Sqd(g_0g_1)^2$.
Therefore $T_u-1\in M\mathcal O_E\cap\mathbb Z=M\mathbb Z$.
For the last equality, $(T_u-1)/M$ is a rational algebraic
integer, hence an integer.  Every divisor of $M$ is consequently
coprime to $T_u$.  Since $16\mid Q_S$, also $T_u\equiv1\pmod{4d}$.
Quadratic reciprocity gives $(-d/T_u)=1$: the $-1$ sign is $1$
and $T_u\equiv1\pmod d$ makes the other symbol $1$.
The literal host prime classification then says $T_u$ is split.
The prime $2$ is the original ramified-two class, so is not in
$\mathcal P$ either.  This is the same norm/congruence datum
used in RSL8 and the same arithmetic host of AH9--AH12.

The \texttt{FamP.drop\_iff} clause now simplifies, on
$\mathcal P$, to
\[
             p\in L.\mathrm{drop}\ \Longleftrightarrow\
                         \mathrm{NoGain}(X_G,\mathcal F.GK,C_r,p),
                                                               \tag{LFS.30}
\]
because its additional exclusions $p\ne T_u$ have just been
verified.  The definition of \texttt{NoGain} lists: a small
order, a narrow power band, a candidate with large share,
a prime above $X''$, a non-inert class, or a nonzero odd-order
Frobenius step in the inert-residue class.
For a first prime the large-share and high-prime branches are
excluded by its definition, and the non-inert branch is
excluded by (LFS.27).  The small-order branch gives the first
alternative in LFS3, with bound $p<Y$.
The power-band branch gives its second alternative; for an
integer prime $p$, the condition $p\ge k^2x^{26}$ is exactly
$p\ge\lceil k^2x^{26}\rceil_+$, including the endpoint.

Here is the exact order dictionary in the other two cases.
The unramified reciprocity calculation of AH9--AH10 identifies
the dyadic projection of the actual Frobenius step $h_p$ with
$p$ in $(\mathbb Z/2^{e_2+2}\mathbb Z)^\times/\{\pm1\}$.
Projection can only divide its order.  If $h_p$ has odd order,
its projection has order dividing both that odd integer and
$2^{e_2}$; hence the projection is zero and has order $1$.
This is strictly smaller than $x^{\epsilon/4}$ for $x>1$.
Such a first prime is at most $X''$, so it again satisfies
LFS3 with bound $\max(Y,X'')$.  These arguments supply every
branch of the literal \texttt{FirstHyp}.

They also show that \texttt{NoGain} forces $VV=\varnothing$.
The first five listed reasons directly contradict its defining
conditions (a band is exactly the candidate's excluded closed
band).  In the odd-order case, the frame implies
$\operatorname{ord}(h_p)\mid o$; since
$o\le x^{\epsilon/100}<x^{\epsilon/4}$, this contradicts
the candidate's lower-order bound.  Thus an unserved prime has
both gain $g_j(p)=0$ and $\tau_P(p)=0$, as required by LFS5.
Also a served prime must satisfy $p\le X''$ by (LFS.30).

Finally suppose a served top is omitted from $T_p$.
The literal definition of \texttt{TTP} retains all tops at a
signal prime and at $p\ge X'$.  Hence this omission forces
$p<X'$ in the inert-residue class.  The failure of
\texttt{NoGain} gives the candidate conditions and
$\Sigma_p\le\theta$, so $VV$ is exactly the good orbit
representatives.  The omitted top's actual representative
$X_G.bs(\mathrm{stepP}(p),t)$ is therefore not good.
Moreover $p$ lies in the original large-order residue-prime
set.  This proves the exact bad-top inclusion used in LFS8.
The second-class implication is immediate from its defining
branch: it is the same candidate and the same $\Sigma_p>\theta$.
% END LFS12

% BEGIN LFS13
\paragraph{LFS13: install all eleven finite-ledger clauses on the original record.}
\label{claim:LFS13}
Assume the original \texttt{LinkG\_at} at $D,L$, with the same
$X_G,\mathcal F.GK,C,b_o$, and use the parameter record
(LFS.22).  LFS11 provides the unit rule and LFS12 the served,
high, and zero-unserved clauses.  We verify the others directly.

The top window says
$a\le a_t\le a+a/(kx^{13})=a(1+h)$, so in particular $a_t\ge0$.
The literal prime exponents in \texttt{ldAt} are
\[
 I_{j,r,p}=\{s<\log_p^{\mathbb N}((a_t)_+)+2:s\text{ odd}\},
 \qquad r=\operatorname{val}(t).                             \tag{LFS.31}
\]
Their weights are the unchanged
$G.\mathrm{fill}(X_G,b_L,j,\mathrm{partner}(p,t,s))$.
At a top this fill is at least $b'/x$, at a thick place at
least $b$, and elsewhere at least $w''$.
Since $w''=b'/(xn)\le b'/x$ and $w''\le b$, every such
weight is at least $w''$.
The discriminant exponent set is
$[1,\log_d^{\mathbb N}((a_t)_+)+1]$.
Its weight is the original $w_{\rm At}$ at the source's
ramified representative; it is nonnegative and, when that
representative lies in $\Psi$, at least one by GFA10's exact
\texttt{NumN\_at} proof.  Outside $\Psi$ its defining
convention is zero.

At a served prime choose any original coefficient place over it.
For existence, the integral basis of the coefficient field makes
$\mathcal O_{E_0}/p\mathcal O_{E_0}$ a nonzero finite ring of
$p^{[E_0:\mathbb Q]}$ elements.  A maximal proper ideal exists
in this finite ring; its inverse image is a nonzero prime over
$(p)$.  Its residue characteristic is $p$,
and it belongs to the saturated support since $p\le X''$.
The \texttt{ParametersFill} equalities identify its selected
tops with the image defining $T_p$, its $\mathrm{gain}$ with
$g_j(p)$, and its correction with $\tau_P(p)$.
They identify each $\mathrm{wPart}$ with the exact fill in
(LFS.31).  Reindex by the injective top labels and by
$J=\{1,\ldots,m\}$.  The third clause of
\texttt{LinkG\_at} is then precisely
\[
 \sum_j\sum_{r\in T_p}(A_{j,r}(p)-\mathrm{loG}(L,C,b_o,p))
       -m\,\mathrm{flG}(L,C,p)
          \le\sum_jg_j(p)-\tau_P(p).                       \tag{LFS.32}
\]
In particular $T_p\subset T$ is the same actual subset, not a
new subset chosen to make the estimate hold.

For the floor envelope replace $C$ by $P.C=\max(C,0)$;
all factors multiplying $C$ are nonnegative for $a\ge1$ and
$p\ge2$.  The resulting formula is exactly $F(p)$ in LFS6,
because $P.f=2n=D.R.f$ and $P.N=D.R.N(m)$.
For the loss envelope put
\[
 A_1=r_Gb'\log x,\qquad
 A_2=b'/x+\eta+S_c+(w''+t)\log x/\delta_1,
\]
and
\[
 B_p=\Lambda a/(pk)+
       \begin{cases}ap^{-L_1(p)}&p<k^2x^{26},\\a/p^2&p\ge k^2x^{26}.
       \end{cases}
\]
All are nonnegative.  The actual \texttt{loG} is its common
$C(S_c+k+\eta)\log_p(2a)$ term, plus a class-restricted
precision term $CA_1$ or $CA_2+b_o$, plus the class-restricted
orbit term $C\log x\,B_p$.  Replace $C$ by its nonnegative
maximum and remove class restrictions; this can only increase
each term.  The low-prime branch lies below $X'$ by LFS10,
and $p<P.p_0$ is exactly $p<k^2x^{26}$ for integer $p$.
Hence this is exactly
$C_p+\mathbf1_{p<X'}(P_p+O_p)$ with the source's $Q_1,Q_2$,
including the single additive $b_o$ in $Q_2$.

Finally the actual
\[
 L_1(p)=\lceil\log_p(k^2x^{78})\rceil_++1
\]
satisfies $p^{L_1(p)-1}\ge k^2x^{78}=K_0$ by monotonicity
of $p^y$ and the ceiling inequality.  This proves the last
clause and completes the literal \texttt{ParametersHyp}.
The finite equation (LFS.32) uses a supplied local gain estimate;
the present argument reconstructs its exact installation rather
than asserting an independent producer for \texttt{LinkG}.
% END LFS13

% BEGIN LFS14
\paragraph{LFS14: the universal original arithmetic reading of the asked sum.}
\label{claim:LFS14}
For every sufficiently late member of the original family,
every tuple $(e',m',AD,N,T',a',w')$ satisfying the literal
\texttt{Reads}, the heterogeneous equality
$AD=\mathrm{arithOf}(D.host)$, and the original cutoff condition,
we have
\[
 \sum_{t\in T'}\sum_{j\in\mathrm{Fin}(m')}
     \mathrm{ASK}(AD,N,a'_t,t,w'_j)
 =\mathrm{ASK}_{\rm original}(L)
 =(\mathrm{ldAt}).\mathrm{ASK}(d,\mathcal P,J,T).             \tag{LFS.33}
\]
We prove both equalities, including the ramified term.

The equalities in \texttt{Reads} first identify $e'$ with
the actual frame and $m'$ with $m$.  Its top-membership
equivalence and the injective integer labels identify $T'$
bijectively with the actual tops.  Its degree clause then gives
$a'_t=a_t$.  Its weight clause identifies, on the entire
arithmetic half $\Phi$, $w'_j$ with the original
$w_{\rm At}(j+1)$ under the original group equivalence
$E:D.R.GG\longrightarrow (\mathbb Z/2\mathbb Z)^2\times A$.
Explicitly, for a point of $\Phi$, its inverse lies in
$\Psi$; use $D.eTau^{-1}$ on that point to get exactly the
code requested by \texttt{Reads.hw}.  Thus this is equality
of weights at each point, not equality of distributions.
GFA10 supplies $w_{\rm At}\ge1$ there.

On the same arithmetic host AH9--AH11 identify the unramified
Frobenius and its two sign coordinates.  At an inert prime
its second multiplicative sign is $-1$.  Its $s$th power
lies outside $\Psi$ exactly when $s$ is odd.  For odd $s$,
the actual partner is
\[
 (\mathbf1_{\mathrm{signal}},t+s h_p),
\]
and multiplication of the two sign coordinates and the
additive frame coordinate shows that its image is exactly
the source's $\rho(\tau_1(t),\mathrm{Frob}_p^s)$.
The fill identities therefore make its weight exactly
$L.\mathrm{wPart}(j,t,p,s)$.
In additive notation this is the partner appearing in
$\mathrm{ASK}(AD,N,\cdot)$, since $E$ carries products to
sums and powers to multiples.  The host's quadratic sign
identifies its inert predicate with the original signal-or-
inert-residue predicate.  Thus every unramified prime summand
agrees term by term.

At $d$ the actual inertia group is
$\{1,\iota_d\}$ by AH11, with $\iota_d\notin\Psi$.
Every Frobenius-power coset has exactly one representative
outside $\Psi$: use $\mathrm{Frob}_d^s$ if it is outside,
and $\mathrm{Frob}_d^s\iota_d$ otherwise.
This is precisely the source's $g_D(s)$.  It lies outside
$\Psi$, and the partner map $\rho(\tau_1(t),\cdot)$ carries
that complement into $\Psi$.  Thus its resulting partner has
the actual weight at least one from GFA10.  The zero-weight
convention outside $\Psi$ mentioned in LFS13 is never used
for this tail argument.
The arithmetic sum over outside-half inertia representatives
therefore has one term, with the same partner and weight,
and the same coefficient $\log d$.  In particular it has
neither an extra factor $2$ nor a missing inertia average.

The cutoff causes no alteration.  For $W>0$ and
$a_t/W\le N$, every integer $M>N$ has
$\max(a_t/M-W,0)=0$.
This proves vanishing of all terms with exponent $s>N$,
because $p^s>s$, and of all primes $p>N$.
Independently, since the original weights on the partner
half are at least one, all terms with $p^s>(a_t)_+$ vanish.
Hence the infinite original prime-power sums can be cut
at exactly the sets (LFS.31) and the discriminant interval
in LFS13.  Primes above $(a_t)_+$ give zero, so extending
their finite universe to (LFS.27) gives zero extra terms.
The top window ensures $(a_t)_+<B$.

Finally replace $j\in\mathrm{Fin}(m)$ by $j+1\in J$
and use (LFS.26).  These finite bijections prove (LFS.33).
An arbitrary weight outside $\Phi$ is immaterial: every
partner in the defining asked sums belongs to $\Phi$.
The threshold needed here is only
$\max(3,\mathcal F.x_1+1)$, which activates the original
fill identities and GFA10's positivity proof.  In particular
it is chosen before the reading, its cutoff $N$, and all its
auxiliary types and functions.
% END LFS14

% BEGIN LFS15
\paragraph{LFS15: the actual same-record prime-class loss setting.}
\label{claim:LFS15}
Suppose $TT\ne\varnothing$, \texttt{ChBadp\_at} holds,
\texttt{ChW\_at} holds, and the original \texttt{LinkG\_at}
holds on the record considered in LFS11--LFS14.
Then every hypothesis used by LFS7--LFS9 is satisfied with
\[
 \begin{split}
 \mathcal P&=\mathrm{PsAt}(L),\quad
 \mathcal C=\mathrm{cands}(L,\mathcal F.GK),\quad
 \mathcal R=\mathrm{largeOrderResiduePrimes}(L),\\
 J&=\mathrm{lifts}(D),\quad T=\mathrm{TTv}(D),\quad
 \Sigma_p=G.\mathrm{Sig}(X_G,\mathcal F.GK,a_L,b_L,p),\\
 io(p)&=G.\mathrm{ioP}(X_G,\mathcal F.GK,a_L,b_L,p),\quad
 D_{\rm fin}=\mathrm{ldAt}(X_G,\mathcal F.GK,C,b_o).
 \end{split}                                                   \tag{LFS.34}
\]
We verify the share clauses that have not yet been checked.
Each small-prime $\kappa$ is a minimum of $1$ and a
nonnegative expression.  Each large-prime $\kappa$ is the
infimum of a set of nonnegative feasible values, using the
literal \texttt{Extra.kapI\_def}; its infimum is nonnegative,
including the empty-set convention.  Thus $\Sigma_p\ge0$,
as also established in GFA4.  The \texttt{ChW\_at} expression
is exactly (LFS.15), because the candidate sets agree and
$x^{-9}=(x^9)^{-1}$.  It has no hidden expectation remaining.

If $VV(p)$ is empty, $io(p)=0$.  Otherwise its literal value is
\[
             io(p)=2\Sigma_p\max_{v\in VV(p)}S_v.
\]
The second clause of \texttt{LinkG\_at} gives
$S_v\le2mN_va/p$.  The original $N_v$ counts a subset of
the tops, so $N_v\le|TT|=q$; this is just the cardinality
bound for the actual filtered top set.
Multiplication by $2\Sigma_p\ge0$ proves the exact
$io(p)\le2\Sigma_p(2mqa/p)$ of LFS8.

By LFS11 the bad subsets have the exact original cardinalities.
Thus \texttt{ChBadp\_at} is exactly the inequality used in
(LFS.18).  LFS12 supplies the first-class, second-class,
and omitted-top implications.  LFS13 supplies every
finite-ledger clause.  Every prime in $\mathcal P$ is odd
by LFS12, and $q\ge1$ by nonemptiness and (LFS.26).
This proves the full literal \texttt{PrimeClassLossHypotheses}
at (LFS.34), or its source's \texttt{LossSettingFamily}
instance with the harmless external index type \texttt{Unit}.
Neither the selected top set nor the local gains are enlarged
to establish these hypotheses.
% END LFS15

% BEGIN LFS16
\paragraph{LFS16: both original losses, with one constant and one threshold.}
\label{claim:LFS16}
Assume, on the same original family,
\[
 \texttt{ChosenTT}(\mathcal F),\qquad
 \texttt{ChosenBadp}(\mathcal F),\qquad
 \texttt{LinkG}(\mathcal F).                              \tag{LFS.35}
\]
Then the complete literal \texttt{LossesSmallShares} holds:
there are $C_L\ge0$ and $x_L$ such that, for all $x\ge x_L$
and every $i\in\mathcal F.I_x$,
\[
 \texttt{ChW\_at}(\mathcal F.Gx(x,i),\mathcal F.GK)
       \Longrightarrow\texttt{Loss\_at}(\mathcal F.L(x,i),C_L,0).
                                                               \tag{LFS.36}
\]
To prove it, take once and for all the witnesses
$C,x_G,\delta$ of \texttt{LinkG}.
Take $e,K$ from (LFS.23), $C_{\rm ord}=12$, and
\[
                              C_L=3|K|+13.                   \tag{LFS.37}
\]
Choose $x_L$ to be the maximum of the \texttt{ChosenTT} and
\texttt{ChosenBadp} thresholds, $x_G$, LFS10's window
threshold, LFS9's uniform numerical threshold, and LFS14's
reading threshold.  All are chosen before $i$ and before
any arithmetic reading.

Fix $x\ge x_L$ and $i$, and suppose the premise of (LFS.36).
The first chosen-top bound is
$|TT|\ge\delta_1 n/2>0$, so $TT$ is nonempty.
LFS10 installs the original parameter record with
$b_o=b^{\delta(x)}$, LFS15 installs the finite loss setting,
and LFS9 gives (LFS.20).  The cardinality identities of
LFS11 identify its $U$ exactly with $|TT|mb$.
Since $2\log4\le3$, $0\le P.C\le K\le|K|$, and
$\sigma\le1$,
\[
 \mathrm{IOsum}(L)\le U\le C_LU,
 \qquad\mathrm{ASK}_{\rm original}(L)-L.\mathrm{Ga}
                  \le(3P.C+13)U\le C_LU.                 \tag{LFS.38}
\]
For the elementary numerical inequality, $\log2<3/4$
follows from $e^{3/4}>1+3/4+(3/4)^2/2>2$, so
$2\log4=4\log2<3$.

Now take an arbitrary \texttt{Reads} tuple, arbitrary
$AD$ heterogeneously equal to the original host arithmetic,
and any cutoff satisfying its specified inequalities.
LFS14 replaces its asked sum by precisely the middle
quantity in (LFS.38).  LFS11 replaces its finite gain and
prime-logarithmic sum by the original $L.\mathrm{Ga}$ and
$\mathrm{IOsum}(L)$.  The source factor
$L.sc.x^{0\cdot L.sc.\epsilon}$ is exactly $1$.
Thus both clauses, with all their universal quantifiers,
are exactly $\texttt{Loss\_at}(L,C_L,0)$.
This proves (LFS.36), not merely a loss estimate for one
preferred reading or one cutoff.
% END LFS16

% BEGIN LFS17
\paragraph{LFS17: preserve the chosen realization and obtain literal \texttt{LossFS}.}
\label{claim:LFS17}
In addition to (LFS.35), assume the original
$\texttt{ChosenWS}(\mathcal F)$, with its own threshold $x_W$.
Let $C_L,x_L$ be supplied by LFS16 and put
$x_0=\max(x_L,x_W)$.  For any $x\ge x_0$ in
$\mathcal F.\mathrm{Win}$, take the exact witnesses of
\texttt{ChosenWS}:
\[
 i_0\in\mathcal F.I_x,\qquad
 a_v,b_v:\mathbb N\longrightarrow\mathbb Z,
 \qquad i_1\in\mathcal F.I_x.                              \tag{LFS.39}
\]
Their stated windows are retained verbatim:
\[
 \forall r\in\mathcal F.\mathrm{TTn}(x),\quad
 \mathrm{aWin}(\mathcal F.L(x,i_0),a_v(r)),\quad
 (\mathcal F.L(x,i_0)).sc.b'\le b_v(r)\le\mathcal F.bTop(x).
                                                               \tag{LFS.40}
\]
The same $i_1$ realizes them at every one of its actual tops:
\[
 (\mathcal F.L(x,i_1)).aT(t)=a_v(\operatorname{val}(t)),\qquad
 (\mathcal F.L(x,i_1)).bT(t)=b_v(\operatorname{val}(t)).         \tag{LFS.41}
\]
For an arbitrary $i\in\mathcal F.I_x$ satisfying these same
two equalities at all its tops, \texttt{ChosenWS} supplies
$\texttt{ChW\_at}$ for that $i$.  Apply (LFS.36), since
$x\ge x_L$, to obtain \texttt{Loss\_at} with the unchanged
$C_L$ and exponent $0$.
The tuple $C_L,x_0,i_0,a_v,b_v,i_1$, (LFS.40), (LFS.41),
and this universal implication is exactly
\[
 \boxed{\quad
 \texttt{ChosenTT}\ \wedge\ \texttt{ChosenBadp}\ \wedge\
 \texttt{LinkG}\ \wedge\ \texttt{ChosenWS}
       \ \Longrightarrow\ \texttt{LossFS}(\mathcal F,0).
 \quad}                                                       \tag{LFS.42}
\]
HW3 supplies the active preselected-pair form of
\texttt{ChosenWS} from its already stated constructor
hypotheses.  The present implication requires only that
literal predicate and does not select a new pair or realize
an arbitrary admissible pair.  In particular the distinct
\texttt{CoverAll} property is absent from (LFS.42).
The source's later \texttt{coverS\_carry} use, starting from
the specific member $i_1$, belongs to the subsequent
same-family ledger assembly.
% END LFS17

\paragraph{Scope at the endpoint.}
The complete finite implication (LFS.42) closes the original
loss installation task on a supplied compatible family.
It leaves the original local arithmetic producer
\texttt{LinkG}, the chosen-top and chosen-bad inequalities,
and the common chosen-weight realization in their exact
stated upstream scopes.  It neither replaces those inputs
by abstract numerically convenient data nor asserts that
the full human proof of RH or unrestricted Dirichlet GRH
has thereby been completed.
